%! TEX root = main.tex

\begin{exercise}\label{ex:calc-1}
Using the $\varepsilon$--$\delta$ definition, prove that
\[\lim_{x\to -1}(3x+4)=1.\]
\end{exercise}

\begin{exercise}\label{ex:calc-2}
Using the $\varepsilon$--$\delta$ definition, prove that
\[\lim_{x\to 1}(x^2+x)=2.\]
\end{exercise}

\begin{exercise}\label{ex:calc-3}
Let $c\in\R$.  Prove directly from the definition that
\[\lim_{x\to a}c=c\]
for every accumulation point $a$ of the domain.
\end{exercise}

\begin{exercise}\label{ex:calc-4}
Let
\[
 f(x)=\begin{cases}
  x^2, & x\ne 2,\\
  17, & x=2.
 \end{cases}
\]
Prove that $\lim_{x\to2}f(x)=4$ directly from the definition.
\end{exercise}

\begin{exercise}\label{ex:calc-5}
Prove from the negation of the definition that
\[\lim_{x\to0}\frac{1}{x}\]
does not exist as a real number.
\end{exercise}

\begin{exercise}\label{ex:calc-6}
Let $f:D\to\R$ and let $a$ be an accumulation point of $D$.  Suppose that
$\lim_{x\to a}f(x)=L$.  Prove that for every $r>0$, there exists a
$\delta>0$ such that
\[
  x\in D\text{ and }0<\abs{x-a}<\delta
  \quad\Longrightarrow\quad \abs{f(x)}<\abs{L}+r.
\]
\end{exercise}

\begin{exercise}\label{ex:calc-7}
Write pseudocode for a procedure which, given a tolerance $\varepsilon>0$,
returns a valid $\delta$ for the claim
\[\lim_{x\to2}x^2=4.\]
Specify its precondition and postcondition, then prove the postcondition from
the $\varepsilon$--$\delta$ definition.
\end{exercise}

